\newcommand{\CaseATitle}{大学新生的选机备忘：把「开学季」放进清单}
\newcommand{\CaseAIntro}{给正在为「开学季」看电脑的大学新生：在课堂、图书馆与宿舍之间安排学习任务，选机时可以先把任务拆开。围绕「到店体验」，下面整理一份有资料依据的核对清单。收藏后带着自己的软件需求逐项看。}
\newcommand{\CaseACTA}{把你常用的软件和使用场景留在评论区，一起补全选机清单。 把这些核对项加入候选清单，比较时统一配置与版本。\par 产品规格与功能请以实际销售版本及适用条件为准。}
\newcommand{\CaseAFirst}{资料里可以确认的是：\par 机身重量：约 1.33 kg（圆键盘含电池）；约 1.32 kg（方键盘含电池）。实际重量依配置和测量方法而异。\par 适用条件：实际机身重量根据配置、制造工艺、测量方法的不同可能有所差异。\par 怎么用这条信息：把携带相关规格列入学习者的对比选择核对清单。具体体验可结合自己的使用方式现场确认。}
\newcommand{\CaseALast}{准备一份「开学季」会用到的软件和外设清单。涉及软件支持范围、销售配置与使用条件的内容，分别向官方资料或门店确认。把这些核对项加入候选清单，比较时统一配置与版本。}
\newcommand{\CaseBTitle}{从职场办公出发，如何了解HUAWEI MateBook 14（鸿蒙版）}
\newcommand{\CaseBIntro}{经常出差的移动办公者在「职场办公」中选购电脑，首先要明确希望完成什么任务。本次以「到店体验」为目标，把能够从资料确认的信息与需要现场确认的事项分开说明，便于读者带着具体问题继续比较。}
\newcommand{\CaseBCTA}{欢迎带着常用软件清单到店体验，并核对所选配置与销售版本。 把这些核对项加入候选清单，比较时统一配置与版本。\par 产品规格与功能请以实际销售版本及适用条件为准。}
\newcommand{\CaseBFirst}{围绕办公使用者的需求，先看资料中的具体表述：\par 机身重量：约 1.33 kg（圆键盘含电池）；约 1.32 kg（方键盘含电池）。实际重量依配置和测量方法而异。\par 适用条件：实际机身重量根据配置、制造工艺、测量方法的不同可能有所差异。\par 选购时的意义：把携带相关规格列入办公使用者的对比选择核对清单。这是一项比较维度，实际选型仍应结合任务、版本和配置一起判断。}
\newcommand{\CaseBLast}{可以沿着「职场办公」的实际流程，把常用软件、外设和文件任务列成清单，逐项确认。资料中的规格只说明相应项目；应用适配、使用习惯与现场体验需要另行核对。把这些核对项加入候选清单，比较时统一配置与版本。 如需沟通选型，可将这份清单和自己的任务一起带到门店，便于围绕具体问题展开交流。}
\newcommand{\CaseCTitle}{开学选机，先想想每天怎么带}
\newcommand{\CaseCIntro}{开学季在几台笔记本之间来回比较，每天要背去教室和图书馆的重量，可以先列进到店核对清单。}
\newcommand{\CaseCCTA}{有相同使用需求，欢迎私信交流，也可以预约到店核对实际销售配置。}
\newcommand{\CaseCFirst}{HUAWEI MateBook 14（鸿蒙版）现有资料里的机身重量分两种键盘口径：约 1.33 kg（圆键盘含电池）、约 1.32 kg（方键盘含电池）。资料同时说明，实际机身重量会因配置、制造工艺、测量方法不同而有差异，配置也可能因销售版本而异，需以实际销售版本为准。所以对比时先弄清自己看的是圆键盘还是方键盘的数字，再把手里那台的实际配置对上，这个重量差异才有参考意义。以上为本地资料口径，未经在线复核。}
\newcommand{\CaseCLast}{HUAWEI MateBook 14（鸿蒙版）现有资料里的机身重量分两种键盘口径：约 1.33 kg（圆键盘含电池）、约 1.32 kg（方键盘含电池）。资料同时说明，实际机身重量会因配置、制造工艺、测量方法不同而有差异，配置也可能因销售版本而异，需以实际销售版本为准。所以对比时先弄清自己看的是圆键盘还是方键盘的数字，再把手里那台的实际配置对上，这个重量差异才有参考意义。以上为本地资料口径，未经在线复核。}
